\chapter{缓存友好的软件优化}

软件优化不是看到循环就展开，看到数组就 SIMD。现代 CPU 的大部分性能问题来自数据形状和执行形状不匹配：数据不连续、跨步太大、复用太晚、写入互相干扰、分支破坏预测、临时对象制造额外流量。缓存友好的优化从数据布局开始，而不是从指令开始。

\section{先减少无意义的数据移动}

推理引擎中，权重、激活、中间张量和输出都会占用内存带宽。若每个算子都把完整中间结果写回内存，再由下一个算子读出，带宽会被快速耗尽。算子融合、原地操作、复用缓冲区和避免临时拷贝，都是减少数据移动的手段。

LCQI 第一阶段的两层 MLP 可以复用中间缓冲区。隐藏层输出需要保存给输出层，但 ReLU 可以在隐藏层写出时直接完成，不需要单独再扫一遍。这个优化很小，却代表了算子融合的基本思想：生产数据时顺手完成简单后处理。

\section{布局决定访问轨迹}

矩阵按行主序存储时，连续访问一行友好，连续访问一列不友好。量化线性层如果每个输出神经元对应一行权重，那么计算一个输出时可以顺序扫描这一行，与输入向量做点积。若后续要并行多个输出神经元，每个线程读取同一输入，写不同输出，权重访问也保持顺序。

\begin{lstlisting}[language=C++,caption={行主序权重点积的访问形状}]
float dot_row_major(const std::int8_t* weights,
                    const float* input,
                    float scale,
                    std::int32_t input_size) {
    float acc = 0.0F;
    for (std::int32_t i = 0; i < input_size; ++i) {
        acc += static_cast<float>(weights[i]) * scale * input[i];
    }
    return acc;
}
\end{lstlisting}

这个循环的权重和输入都是顺序访问，适合缓存和预取器。后续要优化时，可以把反量化乘 scale 移到更合适的位置，可以使用多个累加器，也可以在支持的硬件上使用低精度点积指令。但基础访问形状要先正确。

\section{分支和边界}

热循环里的分支会影响预测和向量化。边界检查、padding、mask、不同 dtype 的分发，如果全部放进最内层循环，性能会很差。常见做法是把边界路径和主路径分开：主路径处理规则长度和对齐数据，尾部路径处理少量剩余元素。算子开发里，这种结构比在每次迭代里判断更清晰。

\section{内存分配和生命周期}

频繁分配释放内存会污染 benchmark，也会让推理延迟不稳定。推理引擎应尽量在模型加载或请求开始时准备缓冲区，热路径复用内存。张量对象要区分拥有数据和引用数据；临时缓冲区要有明确生命周期；跨线程共享缓冲区要明确所有权。

\begin{warningbox}
不要在算子热路径里隐藏分配。一个看似普通的 \code{std::vector} 临时对象，可能在每次推理中分配内存，导致延迟抖动。教学代码可以简单，但 benchmark 前必须检查热路径分配。
\end{warningbox}

\section{优化顺序}

本册建议的优化顺序是：先保证正确性，再记录测量边界；先减少大块无意义数据移动，再考虑 SIMD；先让单线程访问轨迹清晰，再加多线程；先检查算法和布局，再看微小指令级调整。这个顺序不保证每次最快，但能避免在错误层级上浪费时间。

\begin{exercisebox}
检查 LCQI 推理路径中是否存在每次推理都会重新分配的 vector。把可以复用的临时缓冲区移动到执行器对象里，再比较推理延迟的平均值和波动。
\end{exercisebox}
